\chapter{परिमेय भिन्नें}\label{ch:b1:fractions}

\omterm{def:b1:fractions:field}{परिमेय भिन्न} बहुपदों का भागफल है। इस छोटे अध्याय की केंद्रीय प्रमेय ---
आंशिक भिन्न वियोजन\index{आंशिक भिन्न वियोजन} --- ऐसे किसी भी भागफल को
प्रारंभिक ईंटों $\frac{c}{(X - a)^k}$ के योग में तोड़ देती है। अपने बीजीय
आकर्षण से आगे, वह परिमेय फलनों के समाकलन (\cref{ch:b1:integration}) और कुछ
श्रेणियों के योग (\cref{ch:b1:series}) की मानक मशीन है।

\section{क्षेत्र $K(X)$}

\begin{definition}\label{def:b1:fractions:field}
$K$ ($= \R$ या $\C$) पर \emph{परिमेय भिन्न}\index{परिमेय भिन्न} एक भागफल $F
= \frac{A}{B}$ है, जहाँ $A, B \in K[X]$, $B \neq 0$; दो भागफल $\frac AB$ और
$\frac{A'}{B'}$ तब अभिन्न माने जाते हैं जब $AB' = A'B$। प्रत्येक भिन्न का एक
\emph{लघुकृत रूप} है जिसमें $\gcd(A, B) = 1$, और वह अचरों तक अद्वितीय है।
स्वाभाविक संक्रियाओं के साथ परिमेय भिन्नों का \omterm{def:b1:logic:sets}{समुच्चय} $K(X)$ एक \omterm{def:b1:structures:field}{क्षेत्र} है।

$F$ (लघुकृत रूप में) के \emph{ध्रुव}\index{ध्रुव} $B$ के मूल हैं; किसी ध्रुव
की \emph{कोटि} $B$ के मूल के रूप में उसकी \omterm{def:b1:poly:derivative}{बहुलता} है। $F$ की \emph{घात} $\deg
F = \deg A - \deg B \in \Z \cup \{-\infty\}$ है।
\end{definition}

\begin{example}[ध्रुव, कोटियाँ और घात पढ़ना]\label{ex:b1:fractions:vocab}
मान लीजिए $F = \dfrac{X^3 - X}{X^4 - 2X^3 + X^2}$। दोनों परतों का गुणनखंडन
कीजिए: अंश $X(X-1)(X+1)$, हर $X^2(X-1)^2$; उभयनिष्ठ गुणनखंड $X(X-1)$ काटिए:
\[
F = \frac{X + 1}{X(X - 1)} \quad\text{(लघुकृत रूप)।}
\]
\omterm{def:b1:fractions:field}{ध्रुव}: $0$ और $1$, दोनों \emph{सरल} --- कोटियाँ लघुकृत हर पर पढ़ी जाती हैं,
अतः मूल हर के प्रत्यक्ष द्विक मूल निरर्थक हैं। घात: $\deg F = 1 - 2 = -1$,
जो अनंतस्पर्शी रूप से दिखाई देती है ($x \to \infty$ होने पर $xF(x) \to 1$)।
घात \omterm{def:b1:poly:def}{बहुपद} की घात की भाँति व्यवहार करती है ($\deg FG = \deg F + \deg G$,
$\deg(F + G) \leq \max$), और यह लेखा-नियम नीचे गुणांक की खोज में लगातार
प्रयुक्त होता है: ``$xF(x)$ की सीमा'' वाला हर तर्क वेश बदली हुई घात-गणना है।
\end{example}

\begin{proposition}[पूर्णांक भाग]\label{prop:b1:fractions:integerpart}
प्रत्येक $F = \frac AB \in K(X)$ अद्वितीय रूप से $F = E + \frac RB$ है, जहाँ
$E \in K[X]$ ($F$ का \emph{पूर्णांक भाग}, अर्थात् \omterm{def:b1:poly:def}{बहुपद} भाग) और $\deg R <
\deg B$। $E \neq 0$ तभी जब $\deg F \geq 0$।
\end{proposition}

\begin{proof}
यूक्लिडीय भाग $A = BE + R$ (\cref{thm:b1:poly:division}) को $B$ से भाग
दीजिए। अद्वितीयता: यदि $E + \frac RB = E' + \frac{R'}{B}$, तो $\deg(R' - R)
< \deg B$ के साथ $(E - E')B = R' - R$, जो $E = E'$ और फिर $R = R'$ पर बाध्य
कर देता है।
\end{proof}

\begin{example}[पहले लघु कीजिए, फिर भाग दीजिए]\label{ex:b1:fractions:integerpart}
$F = \dfrac{X^3 + 1}{X^2 - 1}$ का पूर्णांक भाग खोजिए। आँख मूँदकर भाग देने
पर: $X^3 + 1 = (X^2 - 1)X + (X + 1)$, अतः $F = X + \frac{X + 1}{X^2 - 1}$।
पर भिन्न लघुकृत नहीं थी: $X^3 + 1 = (X + 1)(X^2 - X + 1)$ और $X^2 - 1 =
(X+1)(X-1)$ में गुणनखंड $X + 1$ साझा है, और
\[
F = \frac{X^2 - X + 1}{X - 1} = X + \frac{1}{X - 1} :
\]
पूर्णांक भाग वही $X$ है, पर भिन्नात्मक भाग सिमटकर एक ही ईंट रह जाता है, और
$-1$ पर का ``\omterm{def:b1:fractions:field}{ध्रुव}'' कभी \omterm{def:b1:fractions:field}{ध्रुव} था ही नहीं। \omterm{def:b1:fractions:field}{ध्रुव} खोजने से पहले सदा न्यूनतम
पदों तक लघु कीजिए: $F$ के \omterm{def:b1:fractions:field}{ध्रुव} \emph{लघुकृत} हर के मूल हैं। (पूर्णांक भाग
सरलीकरण से अप्रभावित रहता है, जैसा \cref{prop:b1:fractions:integerpart} की
अद्वितीयता गारंटी देती है।)
\end{example}

\section{$\C$ पर आंशिक भिन्न वियोजन}

\begin{theorem}[$\C$ पर वियोजन]\label{thm:b1:fractions:complex}
मान लीजिए $F = \frac AB \in \C(X)$ लघुकृत रूप में है, जहाँ $B = c\,(X -
a_1)^{m_1} \cdots (X - a_r)^{m_r}$। तब $F$ अद्वितीय रूप से यह है
\[
F = E + \sum_{i=1}^{r} \sum_{k=1}^{m_i}
\frac{c_{i,k}}{(X - a_i)^k},
\qquad E \in \C[X],\ c_{i,k} \in \C ,
\]
जहाँ $E$ $F$ का पूर्णांक भाग है।
\end{theorem}

\begin{proof}
\cref{prop:b1:fractions:integerpart} से हम $\deg A < \deg B$ मान सकते हैं और
$E = 0$ के साथ योग-वियोजन सिद्ध कर सकते हैं।

\emph{ध्रुवों को अलग करना।} $B_1(a_1) \neq 0$ के साथ $B = (X - a_1)^{m_1}
B_1$ लिखिए। \omterm{def:b1:poly:def}{बहुपद} $(X-a_1)^{m_1}$ और $B_1$ \omterm{cor:b1:arith:bezout}{सहअभाज्य} हैं (कोई उभयनिष्ठ मूल
नहीं), अतः $\C[X]$ में बेज़ू से (\cref{ch:b1:poly} की टिप्पणी देखिए) ऐसे
$U_0, V_0$ हैं कि $U_0 (X-a_1)^{m_1} + V_0 B_1 = 1$; $A$ से गुणा करके और $U
= AU_0$, $V = AV_0$ रखकर, फिर $B$ से भाग देकर:
\[
\frac AB = \frac{V}{(X - a_1)^{m_1}} + \frac{U}{B_1} .
\]
घात की शर्तें लागू की जा सकती हैं: $V$ को $(X-a_1)^{m_1}$ से भाग दीजिए, मान
लीजिए $\deg V_1 < m_1$ के साथ $V = (X-a_1)^{m_1}Q + V_1$, और भागफल को दूसरे
पद में समेट लीजिए ($U_1 = U + QB_1$):
\[
\frac AB = \frac{V_1}{(X - a_1)^{m_1}} + \frac{U_1}{B_1},
\qquad \deg V_1 < m_1 ;
\]
घातों की तुलना ($\deg A < \deg B$ और $\deg V_1 < m_1$) से $\deg U_1 < \deg
B_1$ भी अनिवार्य हो जाता है। $\frac{U_1}{B_1}$ पर, ध्रुव-दर-ध्रुव, दोहराने
से सब कुछ नीचे दी गई \omterm{def:b1:fractions:field}{एक-ध्रुव} वाली स्थिति पर आ जाता है।

\emph{एक \omterm{def:b1:fractions:field}{ध्रुव}।} $\deg A < m$ के साथ $\frac{A}{(X-a)^m}$ के लिए: $A$ का $(X
- a)$ की घातों में प्रसार कीजिए, $A = \sum_{j=0}^{m-1} \alpha_j (X - a)^j$
(\omterm{def:b1:poly:def}{बहुपद} का टेलर प्रसार, जैसा \cref{prop:b1:poly:multiplicity} की उपपत्ति में
है); भाग देने पर ठीक ईंटें $\frac{\alpha_j}{(X-a)^{m-j}}$ मिल जाती हैं।
मूर्त रूप में, $\frac{X^2 + 1}{(X+2)^3}$ के लिए: $X = Y - 2$ प्रतिस्थापित
करने पर,
\[
X^2 + 1 = (Y - 2)^2 + 1 = Y^2 - 4Y + 5 ,
\qquad\text{अतः}\qquad
\frac{X^2+1}{(X+2)^3} = \frac1{Y} - \frac4{Y^2} + \frac5{Y^3}
\]
$Y = X + 2$ के साथ: तीनों ईंटें केवल विस्थापित प्रसार को भाग देने से प्रकट
हो जाती हैं --- जब भी एक ही उच्च-कोटि \omterm{def:b1:fractions:field}{ध्रुव} हो तब यही सबसे तेज़ मार्ग है, और
\cref{exo:b1:fractions:3} के लिए यही अनुशंसित भी है।

\emph{अद्वितीयता।} मान लीजिए दो वियोजन संपाती हैं; उनका अंतर एक सर्वसमिका $0
= \sum_{i,k} \frac{d_{i,k}}{(X - a_i)^k}$ है। पूरे को $(X - a_1)^{m_1}$ से
गुणा कीजिए: $i = 1$, $k = m_1$ वाले पद को \emph{छोड़कर} हर पद में $a_1$ पर
लुप्त होने वाला गुणनखंड आ जाता है, और उस पद का गुणांक $d_{1,m_1}$ तथा कम से
कम एक गुणनखंड $(X - a_1)$ लिए हुए पद बन जाता है। $a_1$ पर मान लेने से (यह
वैध है: गुणा के बाद $a_1$ पर कोई \omterm{def:b1:fractions:field}{ध्रुव} नहीं बचता) $d_{1,m_1} = 0$ मिलता है।
शीर्ष गुणांक हट जाने पर $(X - a_1)^{m_1 - 1}$ के साथ दोहराइए, और इसी प्रकार
नीचे $k = 1$ तक; फिर अगले \omterm{def:b1:fractions:field}{ध्रुव} पर जाइए। सभी $d_{i,k}$ लुप्त हो जाते हैं:
अतः वियोजन अद्वितीय है।
\end{proof}

\begin{method}[गुणांकों की संगणना]\label{met:b1:fractions:compute}
व्यवहार में बेज़ू से बचिए; इन्हें मिलाइए:
\begin{enumerate}
  \item \emph{उच्चतम घात के लिए आवरण विधि}: $\frac{1}{(X-a)^m}$ का गुणांक
  ($m$ \omterm{def:b1:fractions:field}{ध्रुव} $a$ की कोटि) है
        \[
        c_{a,m} = \Bigl[\,(X - a)^m F\,\Bigr]_{X = a};
        \]
        $F = \frac AB$ के \emph{सरल} \omterm{def:b1:fractions:field}{ध्रुव} के लिए यह $\frac{A(a)}{B'(a)}$
        है;
  \item शेष गुणांकों के लिए रैखिक संबंध जुटाने हेतु सुविधाजनक बिंदुओं पर
  \emph{मान} और $x \to \infty$ होने पर $xF(x)$ की \emph{सीमाएँ};
  \item जहाँ उपलब्ध हों वहाँ काम आधा करने के लिए \emph{सम-विषमता} या
  संयुग्मन की सममितियाँ।
\end{enumerate}
\end{method}

\begin{remark}[आंशिक भिन्नों की सामान्य भूलें]\label{rem:b1:fractions:pitfalls}
\begin{enumerate}
  \item \emph{पूर्णांक भाग छोड़ देना।} ईंट-वियोजन घात $< 0$ वाली भिन्नों पर
  लागू होता है; $\deg A \geq \deg B$ होने पर पहले भाग दीजिए
  (\cref{prop:b1:fractions:integerpart}), अन्यथा गुणांक की खोज विरोधाभास
  पैदा कर देगी।
  \item \emph{आवरण विधि का सीमा से बाहर प्रयोग।} $(X - a)^k$ से गुणा करके
  $a$ पर मान लेने से गुणांक केवल $k = m$ के लिए मिलता है, अर्थात् \omterm{def:b1:fractions:field}{ध्रुव} की
  \emph{पूरी} कोटि के लिए; निम्न कोटि के गुणांकों के लिए दूसरे संबंध
  (सीमाएँ, मान) चाहिए --- \cref{ex:b1:fractions:multiple} देखिए।
  \item \emph{लघु करना भूल जाना।} \omterm{def:b1:fractions:field}{ध्रुव} \emph{लघुकृत} रूप पर पढ़े जाते हैं;
  अंश और हर का उभयनिष्ठ गुणनखंड भूतिया \omterm{def:b1:fractions:field}{ध्रुव} बना देता है
  (\cref{ex:b1:fractions:integerpart})।
  \item \emph{$\R$ पर ईंटों का ग़लत आकार।} किसी अखंडनीय द्विघात के ऊपर अंश
  \emph{ऐफ़ीन} ($\alpha X + \beta$) होते हैं, अचर नहीं; केवल $\frac{c}{X^2 +
  1}$ लिखने से हल खो जाते हैं --- सही आकार \cref{thm:b1:fractions:real} तय
  करती है, उन्हें कभी अपने मन से मत गढ़िए।
\end{enumerate}
\end{remark}

\begin{proof}[सरल-ध्रुव सूत्र की उपपत्ति]
सरल \omterm{def:b1:fractions:field}{ध्रुव} $a$ के पास: $Q(a) \neq 0$ के साथ $B = (X - a) Q$, और $B' = Q + (X
- a) Q'$, अतः $B'(a) = Q(a)$। आवरण-मान $\frac{A(a)}{Q(a)} =
\frac{A(a)}{B'(a)}$ है।
\end{proof}

\begin{example}\label{ex:b1:fractions:simple}
$F = \dfrac{1}{X(X-1)(X-2)}$ का वियोजन कीजिए। तीन सरल \omterm{def:b1:fractions:field}{ध्रुव}; हर एक पर आवरण
विधि:
\[
c_0 = \frac{1}{(0-1)(0-2)} = \frac12,
\quad
c_1 = \frac{1}{1 \times (1 - 2)} = -1,
\quad
c_2 = \frac{1}{2 \times 1} = \frac12,
\]
अतः $F = \dfrac{1/2}{X} - \dfrac{1}{X-1} + \dfrac{1/2}{X-2}$। \emph{$X = 3$
पर जाँच:} सीधे $F(3) = \frac{1}{3\cdot2\cdot1} = \frac16$; वियोजन से
$\frac{1/2}{3} - \frac{1}{2} + \frac{1/2}{1} = \frac16 - \frac12 + \frac12 =
\frac16$।
\end{example}

\begin{example}[बहु ध्रुव]\label{ex:b1:fractions:multiple}
$F = \dfrac{X}{(X-1)^2 (X+1)}$ का वियोजन कीजिए। आकार: $\frac{a}{(X-1)^2} +
\frac{b}{X-1} + \frac{c}{X+1}$।
\begin{itemize}
  \item द्विक \omterm{def:b1:fractions:field}{ध्रुव} पर आवरण विधि: $a = \bigl[\frac{X}{X+1}\bigr]_{X=1} =
  \frac12$।
  \item $-1$ पर आवरण विधि: $c = \bigl[\frac{X}{(X-1)^2}\bigr]_{X=-1} =
  -\frac14$।
  \item $\infty$ पर $xF(x)$ की सीमा: $0 = b + c$, अतः $b = \frac14$।
\end{itemize}
\[
F = \frac{1/2}{(X-1)^2} + \frac{1/4}{X-1} - \frac{1/4}{X+1} .
\]
\emph{$X = 0$ पर जाँच:} $F(0) = 0$ और $\frac12 - \frac14 - \frac14 = 0$।
\end{example}

\section{$\R$ पर वियोजन}

\begin{theorem}[$\R$ पर वियोजन]\label{thm:b1:fractions:real}
मान लीजिए $F \in \R(X)$ लघुकृत रूप में है और उसका हर है
\[
B = c \prod_i (X - a_i)^{m_i} \prod_j (X^2 + p_j X + q_j)^{n_j}
\qquad (p_j^2 - 4q_j < 0).
\]
तब $F$ अद्वितीय रूप से अपने पूर्णांक भाग और इन पदों में वियोजित होता है
\[
\frac{c_{i,k}}{(X - a_i)^k} \quad (1 \leq k \leq m_i),
\qquad
\frac{\alpha_{j,l}\, X + \beta_{j,l}}{(X^2 + p_j X + q_j)^{l}}
\quad (1 \leq l \leq n_j),
\]
जिनके गुणांक वास्तविक हैं।
\end{theorem}

\begin{proof}
$\C$ (\cref{thm:b1:fractions:complex}) पर वियोजन कीजिए। चूँकि $F$ वास्तविक
है, \omterm{def:b1:fractions:field}{ध्रुव} $\conj a$ के ऊपर का गुणांक (हर कोटि पर) $a$ के ऊपर के गुणांक का
संयुग्मी है (वियोजन पर संयुग्मन लगाइए और अद्वितीयता का आह्वान कीजिए)। हर
संयुग्मी जोड़े को समूहबद्ध कीजिए:
\[
\frac{c}{(X - z)^l} + \frac{\conj c}{(X - \conj z)^l}
= \frac{c\,(X - \conj z)^l + \conj c\,(X - z)^l}{\bigl(X^2 - 2\Re(z)X
+ \abs z^2\bigr)^{l}},
\]
जिसका अंश अपना ही संयुग्मी है, अतः वास्तविक है, और उसकी घात $\leq l$ है;
वास्तविक द्विघात के गुणजों को अलग करने पर वह हर स्तर $l$ पर घात $\leq 1$ तक
गिर जाता है (एक छोटा अवरोही आगमन)। वास्तविक \omterm{def:b1:fractions:field}{ध्रुव} अपने वास्तविक गुणांक बनाए
रखते हैं (संयुग्मन उन्हें अचल रखता है)। अद्वितीयता $\C$ पर की अद्वितीयता से
निकलती है।
\end{proof}

\begin{example}[संयुग्मी युग्मन को देखना]\label{ex:b1:fractions:pairing}
उपपत्ति की क्रियाविधि, सबसे छोटी स्थिति पर: $\C$ पर $\frac1{X^2+1}$ के \omterm{def:b1:fractions:field}{ध्रुव}
$\pm\iu$ हैं, जिनके आवरण-गुणांक $\iu$ पर $\frac1{2\iu}$ और $-\iu$ पर
$\frac1{-2\iu}$ हैं --- जो एक-दूसरे के संयुग्मी हैं, जैसा प्रमेय बताती है:
\[
\frac{1}{X^2 + 1}
= \frac{1/(2\iu)}{X - \iu} - \frac{1/(2\iu)}{X + \iu} .
\]
उभयनिष्ठ हर पर फिर से जोड़ने पर:
\[
\frac{1}{2\iu}\cdot\frac{(X + \iu) - (X - \iu)}{X^2 + 1}
= \frac{1}{2\iu}\cdot\frac{2\iu}{X^2 + 1} = \frac{1}{X^2+1} :
\]
काल्पनिक भाग कट जाते हैं और वास्तविक ईंट ज्यों-की-त्यों लौट आती है। वास्तविक
समाकल्यों के लिए प्रायः $\R$ छोड़ना ही नहीं पड़ता --- पर जब सम्मिश्र बिंदुओं
पर \emph{योग} निकाले जाते हैं (जैसा सप्ताहांत समस्या इकाई के मूलों के साथ
करती है), तब सम्मिश्र ईंटें ही स्वाभाविक मुद्रा हैं, और यह युग्मन दोनों
वियोजनों के बीच विनिमय-दर है।
\end{example}

\begin{example}\label{ex:b1:fractions:real}
$\R$ पर $F = \dfrac{4}{(X^2+1)(X-1)^2}$ का वियोजन कीजिए। आकार: $\frac{aX +
b}{X^2 + 1} + \frac{c}{(X-1)^2} + \frac{d}{X - 1}$। द्विक \omterm{def:b1:fractions:field}{ध्रुव} पर आवरण
विधि: $c = \bigl[\frac{4}{X^2+1}\bigr]_{X=1} = 2$। सम्मिश्र \omterm{def:b1:fractions:field}{ध्रुव} $\iu$ पर
आवरण विधि ($X^2 + 1$ के ऊपर का अंश सम्मिश्र वियोजन से, या सीधे): $X^2 + 1$
से गुणा कीजिए और $X = \iu$ रखिए:
\[
a\iu + b = \frac{4}{(\iu - 1)^2} = \frac{4}{-2\iu} = 2\iu ,
\]
अतः $a = 2$, $b = 0$। अनंत पर $xF(x)$ की सीमा: $0 = a + d$, अतः $d = -2$।
अतः
\[
F = \frac{2X}{X^2+1} + \frac{2}{(X-1)^2} - \frac{2}{X-1} .
\]
\emph{$X = 0$ पर जाँच:} $F(0) = 4$ और $0 + 2 + 2 = 4$।
\end{example}

\begin{example}[त्रिघात हर, आरंभ से अंत तक]\label{ex:b1:fractions:cubic}
$\R$ पर $F = \dfrac{1}{X^3 + 1}$ का वियोजन कीजिए। पहले गुणनखंडन: $X^3 + 1 =
(X + 1)(X^2 - X + 1)$, जहाँ द्विघात का विविक्तकर $-3 < 0$ है। आकार:
$\frac{a}{X+1} + \frac{bX + c}{X^2 - X + 1}$। सरल \omterm{def:b1:fractions:field}{ध्रुव} $-1$ पर आवरण विधि:
$a = \bigl[\frac1{X^2 - X + 1}\bigr]_{X=-1} = \frac13$। अनंत पर $xF(x)$ की
सीमा: $0 = a + b$, अतः $b = -\frac13$। $X = 0$ पर मान: $1 = a + c$, अतः $c =
\frac23$। अतः
\[
\frac{1}{X^3+1} = \frac13\Bigl(\frac{1}{X+1}
+ \frac{-X + 2}{X^2 - X + 1}\Bigr) ,
\]
जिसकी पुष्टि $X = 1$ पर हो जाती है: बायाँ पक्ष $\frac12$, दायाँ पक्ष
$\frac13\bigl(\frac12 + 1\bigr) = \frac12$। मितव्ययिता पर ध्यान दीजिए: तीन
अज्ञात, तीन सस्ते रैखिक तथ्य (एक आवरण, एक सीमा, एक मान), और किसी चीज़ का कोई
प्रसार नहीं --- अर्थात् \cref{met:b1:fractions:compute} का कार्यप्रवाह अपने
शुद्ध रूप में।
\end{example}

\begin{remark}[यह किस काम आता है]
वियोजन हो जाने पर परिमेय फलन पद-दर-पद समाकलित हो जाता है: ईंटों
$\frac{1}{(x-a)^k}$ के प्रारंभिक प्रतिअवकलज हैं, और ईंटें $\frac{\alpha x +
\beta}{(x^2 + px + q)^l}$ $\ln$ तथा $\arctan$ तक सिमट जाती हैं
(\cref{ch:b1:integration})। दूरबीनी योग दूसरा मानक अनुप्रयोग हैं
(\cref{exo:b1:fractions:8})।
\end{remark}

\begin{remark}[यह अध्याय कहाँ काम आता है]\label{rem:b1:fractions:whereused}
आंशिक भिन्नें सबसे बढ़कर एक \emph{पूर्वसंस्करण} पद हैं: समाकलन का अध्याय
(\cref{ch:b1:integration}) हर परिमेय समाकल्य को समाकलन से पहले
\cref{thm:b1:fractions:real} से गुज़ारता है, और श्रेणी का अध्याय
(\cref{ch:b1:series}) परिमेय पदों को ठीक वैसे ही दूरबीनी करता है जैसे
\cref{exo:b1:fractions:8} में और नीचे दी गई सप्ताहांत समस्या में --- जो इस
तकनीक को $\sum 1/k^2 = \pi^2/6$ तक ले जाती है। लघुगणकीय अवकलज $P'/P = \sum
m_i/(X - a_i)$ (\cref{exo:b1:poly:11}) जब भी मूलों के स्थान का अध्ययन होता
है तब फिर प्रकट होता है। इस खंड से आगे, \omterm{def:b1:diffeq:linear2}{अभिलक्षणिक बहुपद} $\chi$ के लिए
$1/\chi(X)$ का वियोजन स्नातक वर्ष 2 के खंड में आव्यूह की घातों और लाप्लास
रूपांतरों की संगणना के मूल में है: ईंटें $\frac{c}{(X - a)^k}$
\cref{ch:b1:diffeq} में मिले हलों $t^{k-1}\eu^{at}$ की बीजीय छाया हैं।
\end{remark}

\begin{omfigure}
\begin{tikzpicture}[xscale=1.75, yscale=0.42]
  \draw[->, thin] (-2.7,0) -- (2.8,0) node[below] {$x$};
  \draw[->, thin] (0,-4.6) -- (0,4.9) node[left] {$y$};
  \draw[dashed, gray] (-1,-4.4) -- (-1,4.6);
  \draw[dashed, gray] (1,-4.4) -- (1,4.6);
  \draw[omThm, thin] (-2.6,1.5) -- (2.7,1.5)
    node[above right, font=\small, pos=0.02] {$y = \lambda$};
  \draw[omDef, thick, domain=-2.6:-1.31, samples=80, smooth]
    plot (\x, {1/(\x+1) + 1/\x + 1/(\x-1)});
  \draw[omDef, thick, domain=-0.83:-0.20, samples=120, smooth]
    plot (\x, {1/(\x+1) + 1/\x + 1/(\x-1)});
  \draw[omDef, thick, domain=0.20:0.83, samples=120, smooth]
    plot (\x, {1/(\x+1) + 1/\x + 1/(\x-1)});
  \draw[omDef, thick, domain=1.31:2.7, samples=80, smooth]
    plot (\x, {1/(\x+1) + 1/\x + 1/(\x-1)});
  \fill[omThm] (-0.715,1.5) circle (0.045 and 0.19);
  \fill[omThm] (0.40,1.5) circle (0.045 and 0.19);
  \fill[omThm] (2.3,1.5) circle (0.045 and 0.19);
\end{tikzpicture}

{\small \cref{exo:b1:fractions:12}-प्रकार का फलन $F(x) = \frac1{x+1} +
\frac1x + \frac1{x-1}$: वह अपने ध्रुवों $-1, 0, 1$ (बिंदुकित) के बीच के
प्रत्येक अंतराल पर कठोरतः ह्रासमान है। पहले \omterm{def:b1:fractions:field}{ध्रुव} के दाईं ओर हर क्षैतिज स्तर
$\lambda > 0$ प्रत्येक अंतराल में ठीक एक बार काटा जाता है (चिह्नित बिंदु):
अर्थात् $F = \lambda$ के हल ध्रुवों के बीच-बीच में आते हैं।}
\end{omfigure}

\section{अभ्यास}

\begin{exercise}[$\star$]\label{exo:b1:fractions:1}
$\R$ पर वियोजन कीजिए: $\dfrac{1}{X^2 - 1}$; $\;\dfrac{X}{X^2 - 3X + 2}$;
$\;\dfrac{X^2 + 1}{X(X-1)}$ \emph{(पूर्णांक भाग का ध्यान रखिए)}।
\end{exercise}

\begin{exercise}[$\star$]\label{exo:b1:fractions:2}
$\R$ पर वियोजन कीजिए: $\dfrac{1}{X(X^2 + 1)}$ और $\dfrac{X^3}{X^2 + X + 1}$।
\end{exercise}

\begin{exercise}[$\star$]\label{exo:b1:fractions:3}
$\dfrac{1}{X^2(X - 1)}$ और $\dfrac{X + 1}{(X - 1)^3}$ का वियोजन कीजिए
\emph{(दूसरे के लिए $Y = X - 1$ प्रतिस्थापित कीजिए)}।
\end{exercise}

\begin{exercise}[$\star\star$]\label{exo:b1:fractions:4}
$\C$ पर, फिर $\R$ पर वियोजन कीजिए: $\dfrac{1}{X^4 - 1}$।
\end{exercise}

\begin{exercise}[$\star\star$]\label{exo:b1:fractions:5}
$\R$ पर वियोजन कीजिए: $\dfrac{X^2}{(X^2 + 1)^2}$, और $\int \frac{\dd
x}{(x^2+1)^2} = \frac12\bigl(\arctan x + \frac{x}{x^2+1}\bigr) + C$ दिया
होने पर $x \mapsto \dfrac{x^2}{(x^2+1)^2}$ का एक प्रतिअवकलज निकालिए।
\end{exercise}

\begin{exercise}[$\star\star$]\label{exo:b1:fractions:6}
$n \in \N^*$ के लिए $F_n = \dfrac{n!}{X(X+1)\cdots(X+n)}$ का वियोजन कीजिए
\emph{($0, -1, \dots, -n$ पर सरल \omterm{def:b1:fractions:field}{ध्रुव}; आवरण सूत्र का प्रयोग कीजिए और \omterm{def:b1:counting:objects}{द्विपद
गुणांक} पहचानिए)}।
\end{exercise}

\begin{exercise}[$\star\star$]\label{exo:b1:fractions:7}
$P = X^n - 1$ के लिए \cref{exo:b1:poly:11} की सर्वसमिका का प्रयोग करके सिद्ध
कीजिए कि
\[
\sum_{k=0}^{n-1} \frac{1}{X - \omega^k} = \frac{n X^{n-1}}{X^n - 1},
\qquad \omega = \eu^{2\iu\pi/n},
\]
और जाँच के लिए $n = 4$ के लिए दोनों पक्षों का मान $X = 2$ पर निकालिए।
\end{exercise}

\begin{exercise}[$\star\star$]\label{exo:b1:fractions:8}
$\dfrac{1}{k(k+1)(k+2)}$ का वियोजन कीजिए और संगणना कीजिए
\[
S_n = \sum_{k=1}^{n} \frac{1}{k(k+1)(k+2)},
\qquad\text{फिर}\qquad
\lim_{n \to \infty} S_n .
\]
\end{exercise}

\begin{exercise}[$\star\star\star$]\label{exo:b1:fractions:9}
मान लीजिए $P \in \R[X]$ घात $n$ का \omterm{def:b1:poly:def}{इकाई-अग्र बहुपद} है जिसके $n$ भिन्न
वास्तविक मूल $x_1 < \dots < x_n$ हैं। सिद्ध कीजिए कि
\[
\sum_{i=1}^{n} \frac{1}{P'(x_i)} = 0 \quad (n \geq 2),
\qquad
\sum_{i=1}^{n} \frac{x_i^{\,n-1}}{P'(x_i)} = 1 .
\]
\emph{संकेत: $m \leq n - 1$ के लिए $\frac{X^m}{P}$ का वियोजन कीजिए और अनंत
पर गुणांकों के क्षय पर दृष्टि डालिए --- अथवा बिंदुओं $x_i$ पर एकपद $X^m$ का
\omterm{thm:b1:poly:lagrange}{लाग्रांज अंतर्वेशन} (\cref{thm:b1:poly:lagrange}) प्रयोग कीजिए।}
\end{exercise}

\begin{exercise}[$\star\star$]\label{exo:b1:fractions:10}
$\R$ पर $\dfrac{1}{X(X+1)^2}$ का वियोजन कीजिए। मान $\sum_{k \geq 1}
\frac1{k^2} = \frac{\pi^2}6$ मानकर (जो इस अध्याय की सप्ताहांत समस्या में
सिद्ध है) निकालिए
\[
\sum_{n=1}^{\infty} \frac{1}{n(n+1)^2} = 2 - \frac{\pi^2}{6} .
\]
\end{exercise}

\begin{exercise}[$\star\star$]\label{exo:b1:fractions:11}
$\R$ पर वियोजन कीजिए: $\dfrac{1}{(X^2+1)(X^2+4)}$, फिर
$\dfrac{X^2}{(X^2+1)(X^2+4)}$। \emph{संकेत: दोनों हर $X^2$ में \omterm{def:b1:poly:def}{बहुपद} हैं:
पहले $\frac1{(Y+1)(Y+4)}$ का वियोजन कीजिए।}
\end{exercise}

\begin{exercise}[$\star\star\star$]\label{exo:b1:fractions:12}
(अभिलक्षणिक-मान समीकरण) मान लीजिए $F = \sum_{i=1}^{r} \frac{c_i}{X - p_i}$,
जहाँ $p_1 < p_2 < \dots < p_r$ वास्तविक हैं और सभी $c_i > 0$।
\begin{enumerate}
  \item दिखाइए कि $F$ अपने प्रांत के प्रत्येक अंतराल पर कठोरतः ह्रासमान है,
  और $\pm\infty$ पर तथा प्रत्येक \omterm{def:b1:fractions:field}{ध्रुव} के दोनों ओर उसकी सीमाएँ दीजिए।
  \item उससे निकालिए कि प्रत्येक $\lambda > 0$ के लिए समीकरण $F(x) =
  \lambda$ के ठीक $r$ वास्तविक हल हैं, प्रत्येक अंतराल
  $\intoo{p_i}{p_{i+1}}$ में एक और $p_r$ से आगे एक। (ऐसे समीकरण अभिलक्षणिक
  मानों के विक्षोभ को नियंत्रित करते हैं; मध्यवर्ती मान प्रमेय यहाँ उच्चतर
  माध्यमिक स्तर पर प्रयुक्त है और \cref{ch:b1:continuity} में सिद्ध।)
\end{enumerate}
\end{exercise}

\section{समस्या: इंजन के रूप में आंशिक भिन्नें}

\begin{problem}\label{pb:b1:fractions:1}
आंशिक भिन्न वियोजन लेखा-जोखा जैसा लगता है; यह समस्या दिखाती है कि वह एक इंजन
है। भिन्न $\frac1{X(X+1)\cdots(X+k)}$ खिलाने पर वह योगों के पूरे कुल संवृत
रूप में दूरबीनी कर देता है; $\frac1{X^n - 1}$ खिलाने पर वह ऐसी त्रिकोणमितीय
सर्वसमिकाएँ उत्पन्न करता है
\[
\sum_{k=1}^{n-1}\frac{1}{\sin^2\frac{k\pi}{n}} = \frac{n^2-1}{3} ;
\]
और एक पद और आगे बढ़ाने पर वही सर्वसमिका गणित के सबसे प्रसिद्ध सूत्रों में से
एक, ऑयलर का सूत्र, निचोड़ देती है
\[
\sum_{k=1}^{\infty}\frac1{k^2} = \frac{\pi^2}{6}
\]
--- और वह भी केवल इस अध्याय के बीजगणित और विद्यालयी त्रिकोणमिति
से।\index{दूरबीनी}\index{बासेल समस्या} आगे सर्वत्र $\omega =
\eu^{2\iu\pi/n}$; अनुक्रमों की सीमाएँ उच्चतर माध्यमिक स्तर पर प्रयुक्त हैं
(\cref{ch:b1:seq} उन्हें औपचारिक बनाता है)।

\medskip \noindent\textbf{भाग I --- दूरबीन।}
\begin{enumerate}
  \item $\frac1{X(X+1)}$ का वियोजन कीजिए और $\sum_{n=1}^{N} \frac1{n(n+1)}$
  की यथार्थ संगणना कीजिए; निष्कर्ष निकालिए कि योग $1$ की ओर जाता है।
  \item $\frac1{X(X+2)}$ के लिए भी वही: दिखाइए $\sum_{n=1}^{N}
  \frac1{n(n+2)} = \frac12\bigl(\frac32 - \frac1{N+1} - \frac1{N+2}\bigr)
  \to \frac34$। (अंतराल होने पर हर सिरे पर \emph{दो} सीमांत पद बच रहते हैं।)
  \item क्रियाविधि को औपचारिक बनाइए: यदि किसी ऐसे परिमेय $G$ के लिए, जिसका
  $\intco1{+\infty}$ में कोई \omterm{def:b1:fractions:field}{ध्रुव} न हो, $F(X) = G(X) - G(X+1)$ हो, तो
  $\sum_{n=1}^N F(n) = G(1) - G(N+1)$। $F = \frac1{X(X+1)(X+2)}$ के लिए
  साक्षी $G$ प्रस्तुत करके \cref{exo:b1:fractions:8} का मान $\frac14$ फिर से
  प्राप्त कीजिए।
  \item व्यापक क्रमगुणित दूरबीन सिद्ध कीजिए: $k \geq 1$ के लिए,
        \[
        \frac{1}{X(X+1)\cdots(X+k)}
        = \frac1k\biggl(\frac{1}{X(X+1)\cdots(X+k-1)}
        - \frac{1}{(X+1)\cdots(X+k)}\biggr),
        \]
        और उससे निकालिए
        \[
        \sum_{n=1}^{\infty}\frac{1}{n(n+1)\cdots(n+k)}
        = \frac{1}{k \cdot k!} .
        \]
        स्थिति $k = 2$ को प्रश्न 3 के सामने जाँचिए।
  \item \cref{exo:b1:fractions:6} के वियोजन का मान सुचयनित बिंदुओं पर लेकर
  सिद्ध कीजिए
        \[
        \sum_{j=0}^{n}(-1)^j\binom nj\,\frac1{j+1}
        = \frac1{n+1},
        \qquad
        \sum_{j=0}^{n}(-1)^j\binom nj\,\frac1{j+2}
        = \frac1{(n+1)(n+2)} .
        \]
\end{enumerate}

\medskip \noindent\textbf{भाग II --- भिन्न $1/(X^n - 1)$।}
\begin{enumerate}[resume]
  \item आवरण सूत्र से दिखाइए कि
        \[
        \frac{1}{X^n - 1}
        = \frac1n\sum_{k=0}^{n-1}\frac{\omega^k}{X - \omega^k} .
        \]
  \item दो जाँच: $n = 2$ के लिए सूत्र सीधे सत्यापित कीजिए, और दिखाइए कि $n
  \geq 2$ के लिए $n$ गुणांकों का योग लुप्त हो जाता है --- समझाइए कि ऐसा होना
  ही चाहिए ($x \to \infty$ होने पर $xF(x)$ पर विचार कीजिए)।
  \item आवरण विधि से \cref{exo:b1:fractions:7} की सर्वसमिका फिर से
  व्युत्पन्न कीजिए: $\frac{nX^{n-1}}{X^n-1} = \sum_k \frac1{X - \omega^k}$।
  \item संयुग्मी ध्रुवों को समूहबद्ध करके वास्तविक वियोजन सिद्ध कीजिए:
  $\theta_k = \frac{2k\pi}n$ के साथ,
        \[
        \frac{\omega^k}{X - \omega^k} +
        \frac{\omega^{n-k}}{X - \omega^{n-k}}
        = \frac{2\cos\theta_k\,X - 2}
        {X^2 - 2\cos\theta_k\,X + 1} ,
        \]
        और $\frac1{X^n-1}$ का पूरा वास्तविक वियोजन लिखिए ($n$ के विषम और $n$
        के सम होने में भेद कीजिए)।
  \item $n = 4$ पर विशेषीकरण कीजिए और \cref{exo:b1:fractions:4} के सामने
  जाँचिए।
\end{enumerate}

\medskip \noindent\textbf{भाग III --- त्रिकोणमितीय योग, और ऑयलर का
$\pi^2/6$।} मान लीजिए $P = 1 + X + \dots + X^{n-1}$, जिसके मूल $\omega,
\omega^2, \dots, \omega^{n-1}$ हैं (सभी सरल)।
\begin{enumerate}[resume]
  \item $\frac{P'}P = \sum_{k=1}^{n-1}\frac1{X - \omega^k}$
  (\cref{exo:b1:poly:11}) का प्रयोग करके सिद्ध कीजिए
        \[
        \sum_{k=1}^{n-1}\frac{1}{1 - \omega^k} = \frac{n-1}2 .
        \]
  \item $\theta \notin 2\pi\Z$ के लिए $\dfrac1{1 - \eu^{\iu\theta}} =
  \dfrac12 + \dfrac\iu2\,\frac{\cos(\theta/2)}{\sin(\theta/2)}$ सिद्ध कीजिए
  (अर्ध-कोण गुणनखंडन, \cref{met:b1:complex:trig}), और प्रश्न 11 से निकालिए
  कि $\sum_{k=1}^{n-1} \frac{\cos(k\pi/n)}{\sin(k\pi/n)} = 0$ --- जो सममिति
  $k \leftrightarrow n - k$ से भी दिखाई देता है।
  \item प्रश्न 11 की सर्वसमिका का अवकलन करके (अर्थात् $X = 1$ पर
  $\bigl(\frac{P'}P\bigr)' = \frac{P''}P - \bigl(\frac{P'}P\bigr)^2$ का मान
  लेकर) सिद्ध कीजिए
        \[
        \sum_{k=1}^{n-1}\frac{1}{(1 - \omega^k)^2}
        = \frac{(n-1)(5-n)}{12} .
        \]
  \item $\cot t = \frac{\cos t}{\sin t}$ लिखकर प्रश्न 12--13 से दोनों संवृत
  रूप निकालिए
        \[
        \sum_{k=1}^{n-1}\cot^2\frac{k\pi}n = \frac{(n-1)(n-2)}3,
        \qquad
        \sum_{k=1}^{n-1}\frac1{\sin^2\frac{k\pi}n}
        = \frac{n^2-1}3 .
        \]
  \item $n = 3$ और $n = 4$ के लिए दोनों सूत्र हाथ से सत्यापित कीजिए।
  \item $t \in \intoo0{\frac\pi2}$ के लिए असमिकाएँ $\cot t < \frac1t <
  \frac1{\sin t}$ सिद्ध कीजिए ($\sin t < t < \tan t$ से), और उससे $n = 2m +
  1$ तथा $1 \leq k \leq m$ के लिए निकालिए:
        \[
        \cot^2\frac{k\pi}n \;<\; \frac{n^2}{k^2\pi^2}
        \;<\; \frac1{\sin^2\frac{k\pi}n} .
        \]
  \item $k = 1, \dots, m$ के लिए इन असमिकाओं का योग कीजिए (प्रश्न 14 के
  सूत्रों को आधा करने के लिए सममिति $k \leftrightarrow n - k$ का प्रयोग
  कीजिए) और दबाइए:
        \[
        \sum_{k=1}^{\infty}\frac1{k^2} = \frac{\pi^2}6 .
        \]
\end{enumerate}

\medskip \noindent\textbf{भाग IV --- ऊँची ईंटें।}
\begin{enumerate}[resume]
  \item $\frac1{X^2-1}$ के वियोजन का वर्ग करके और आड़े पद का फिर से वियोजन
  करके सिद्ध कीजिए
        \[
        \frac1{(X^2-1)^2}
        = \frac14\Bigl(\frac1{(X-1)^2} + \frac1{(X+1)^2}\Bigr)
        - \frac14\Bigl(\frac1{X-1} - \frac1{X+1}\Bigr),
        \]
        और उसे $X = 0$ पर जाँचिए।
  \item प्रश्न 18, दूरबीन और ऑयलर का मान (प्रश्न 17) मिलाकर सिद्ध कीजिए
        \[
        \sum_{n=2}^{\infty}\frac1{(n^2-1)^2}
        = \frac{\pi^2}{12} - \frac{11}{16} ,
        \]
        और मान की पुष्टि तीन दशमलव तक संख्यात्मक रूप से कीजिए।
  \item प्रश्न 8 की सर्वसमिका का अवकलन करके $\sum_{k=0}^{n-1}\frac1{(X -
  \omega^k)^2}$ का संवृत रूप प्राप्त कीजिए, और उसे $X = 2$, $n = 2$ पर
  जाँचिए।
  \item $k \geq 2$ के लिए सिद्ध कीजिए कि
        \[
        \sum_{n=k}^{\infty}\frac{1}{\binom nk} = \frac{k}{k-1} .
        \]
        \emph{($1/\binom nk$ को क्रमगुणितों से लिखकर प्रश्न 4 पर ले आइए।)}
  \item $k = 3$ के लिए यथार्थ आंशिक योग $\sum_{n=3}^{N}\frac1{\binom n3}$ और
  उसकी सीमा दीजिए।
\end{enumerate}

\medskip \noindent\textbf{भाग V --- संश्लेषण।}
\begin{enumerate}[resume]
  \item शिखर-संगणना के रूप में $\dfrac1{X^6 - 1}$ का पूरा वास्तविक वियोजन
  लिखिए और उसे $X = 0$ पर जाँचिए।
  \item इस समस्या में ठीक कहाँ प्रयोग हुआ: (क) वियोजन की अद्वितीयता; (ख)
  \cref{ch:b1:complex} से \omterm{def:b1:complex:unity}{इकाई के मूल}; (ग) \cref{exo:b1:poly:11} से लघुगणकीय
  अवकलज? प्रत्येक के लिए एक वाक्य।
  \item एक छोटे अनुच्छेद में संश्लेषण: एक ही बीजीय सर्वसमिका --- भिन्न को
  ईंटों में तोड़ना --- ने यथार्थ योग, त्रिकोणमितीय सर्वसमिकाएँ और $\pi^2/6$
  उत्पन्न किए। बीजगणित (यथार्थ वियोजन, सर्वत्र वैध) और विश्लेषण (सीमाएँ,
  दबाना) के बीच श्रम-विभाजन पर टिप्पणी कीजिए, और बताइए कि प्रत्येक धागा कहाँ
  औद्योगिक रूप ले लेता है: \cref{ch:b1:series} में दूरबीनी और तुलना, तथा
  \cref{ch:b1:integration} में ईंटों का समाकलन।
\end{enumerate}
\end{problem}
